\documentclass[a4paper,12pt]{article}
\usepackage[T2A]{fontenc}
\usepackage[utf8]{inputenc}
\usepackage{cmap}
\usepackage[english, russian]{babel}
\usepackage{wrapfig}
\usepackage[backend=biber,style=gost-numeric]{biblatex} % стиль ГОСТ
\addbibresource{literatura.bib}
\usepackage{misccorr} % в заголовках появляется точка, но при ссылке на них ее нет
\usepackage{amssymb,amsfonts,amsmath,amsthm}  
\usepackage{indentfirst}
\usepackage[usenames,dvipsnames]{color} 
\usepackage[unicode,hidelinks]{hyperref}
% \hypersetup{%
%     pdfborder = {0 0 0}
\usepackage{multirow}
\usepackage{makecell,multirow} 
\usepackage{ulem}
%\usepackage{hyperref}
%\usepackage{graphicx,wrapfig}
%\renewcommand\epsilon{\varepsilon}
% \renewcommand\epsilon{\varepsilon}
%\graphicspath{{img/}}
\usepackage{geometry}
\geometry{left=1cm,right=1cm,top=2cm,bottom=2cm,bindingoffset=0cm,headheight=15pt}
\usepackage{fancyhdr} 
\linespread{1.2} 
\frenchspacing 
\renewcommand{\labelenumii}{\theenumii)} 
% \usepackage{caption}
%%%%%%%%%%%%%%%%%%%%%%%%%%%%%%%%%%%%%%%%%%%%%%%%%%%%%%%%%%%%%%%%%%%%%%%%%%%%%%%
%%%%%%%%%%%%%%%%%%%%%%%%%%%%%%%%%%%%%%%%%%%%%%%%%%%%%%%%%%%%%%%%%%%%%%%%%%%%%%%

\def\labauthor{Сарафанов И.Г.}
\def\labauthors{\labauthor}
\def\labnumber{127}
\def\labtheme{Таблица производных и интегралов}

%%%%%%%%%%%%%%%%%%%%%%%%%%%%%%%%%%%%%%%%%%%%%%%%%%%%%%%%%%%%%%%%%%%%%%%%%%%%%%%
	%применим колонтитул к стилю страницы
\pagestyle{fancy} 
	%очистим "шапку" страницы
\fancyhead{\LaTeX} 
	%слева сверху на четных и справа на нечетных
\fancyhead[L]{\labauthors} 
	%справа сверху на четных и слева на нечетных
\fancyhead[R]{Яковлев \textnumero 237} 
	%очистим "подвал" страницы
\fancyfoot{} 
	% номер страницы в нижнем колинтуле в центре
\fancyfoot[C]{\thepage} 
%%%%%%%%%%%%%%%%%%%%%%%%%%%%%%%%%%%%%%%%%%%%%%%%%%%%%%%%%%%%%%%%%%%%%%%%%%%%%%%
\usepackage{float}
\usepackage[mode=buildnew]{standalone}
\usepackage{tikz} 
% \usepackage{subcaption}
\usepackage{tikz,csvsimple,physics}
\makeatother
%\usepackage{booktabs}
\usepackage{pgfplots, pgfplotstable}
\usepackage[outline]{contour}
\usepackage{tocloft}
\renewcommand{\cftsecleader}{\cftdotfill{\cftdotsep}}


% вот этот удивительный мир

\begin{document}


%\textit{Дано: $\textbf{u}$ -- скорость газа относительно ракеты, $\mu$ -- расход газа, показать, что уравнение движения $ma = F - \mu u$ }

\textit{Условие:} На носу лодки длиной $l$ стоит человек, держа на высоте $h$ ядро массы $m$. Масса лодки вместе с человеком равна $M$. Человек бросает горизонтально ядро вдоль лодки. Какую скорость по горизонтали должен сообщить человек ядру, чтобы попасть в корму лодки? Сопротивлением воды движению лодки можно пренебречь.

\vspace{2em}
До броска общий импульс системы равен нулю. После броска ядро движется со скоростью $v$ относительно лодки, лодка движется со скоростью $V$ в противоположном направлении.

\vspace{1em}
Из \textbf{закона сохранения импульса} (в проекции на горизонтальную ось):
\begin{equation}
    m(v - V) = MV
\end{equation}

Отсюда находим скорость лодки:
\begin{equation}
    V = \frac{mv}{M + m}
\end{equation}

Скорость ядра относительно воды:
\begin{equation}
    v_{\text{ядра}} = v - V = v \cdot \frac{M}{M + m}
\end{equation}

Ядро падает с высоты $h$ под действием силы тяжести. Время падения:
\begin{equation}
    t = \sqrt{\frac{2h}{g}}
\end{equation}

За время полёта $t$ ядро должно пролететь расстояние $l$ \textit{относительно лодки}:
\begin{equation}
    (v - V) \cdot t = l
\end{equation}

Подставляем выражения для $V$ и $t$:
\begin{equation}
    v \cdot \frac{M}{M + m} \cdot \sqrt{\frac{2h}{g}} = l
\end{equation} 

Отсюда находим необходимую для попадания в корму лодки скорость:
\begin{equation}
    \boxed{v = \frac{M + m}{M} \cdot l \sqrt{\frac{g}{2h}}}
\end{equation}

\end{document}